\chapter*{Abstract}
\addcontentsline{toc}{chapter}{Abstract}

\begin{itemize}
  \item \textbf{Problem:} animal-movement theory is difficult to apply to real tracking studies because observations, environmental rasters, coordinate systems, sampling intervals, behavioral annotations, and model interfaces rarely align without substantial preprocessing.
  \item \textbf{Objective:} close this theory--practice gap with a complete and reusable MoveApps workflow rather than present conditioned random walks as an isolated theoretical contribution.
  \item \textbf{Workflow:} environmental annotation; metric projection; data-quality checks; trajectory segmentation; behavioral-state annotation; empirical kernel construction; terrain-aware endpoint-conditioned movement reconstruction; utilization distributions; visit probabilities; and lightweight interactive visualization.
  \item \textbf{Methodological basis:} use the discretized conditioned-random-walk framework of \textcite{buchin2024randomwalks}, its environment-aware extensions \parencite{buchin2026advanced}, and the broader account in \textcite{hagedoornPhD}, while introducing every required concept without assuming prior expertise.
  \item \textbf{Software contribution:} expose the individual stages as interoperable MoveApps components with explicit data contracts, user-facing settings, diagnostics, provenance, and recoverable artifacts.
  \item \textbf{Empirical contribution:} evaluate complete workflows on realistic animal-movement data and compare segmentation, kernel, interpolation, and spatial-summary choices against relevant references and ablations.
  \item \textbf{Results to insert:} final quantitative findings from the frozen experiment manifests under \code{/home/omar/Desktop/MSc Thesis/Experiments/} and \code{/home/omar/PycharmProjects/RW-Python-gitlab/benchmark/}.
\end{itemize}

\begin{itemize}
  \item \textbf{Keywords:} animal movement; MoveApps; environmental annotation; trajectory segmentation; movement kernels; conditioned random walks; utilization distribution.
\end{itemize}

\cleardoublepage
